\documentclass[9pt,a4paper,twocolumn]{ctexart}

% --- Packages ---
\usepackage[top=1.2cm, bottom=1.2cm, left=1.2cm, right=1.2cm, columnsep=0.8cm]{geometry}
\usepackage{hyperref}
\usepackage{graphicx}
\usepackage{float}
\usepackage{longtable}
\usepackage{tabularx}
\usepackage{booktabs}
\usepackage{enumitem}
\usepackage{xcolor}
\usepackage{fancyhdr}
\usepackage{listings}
\usepackage{fontspec}
\usepackage{titlesec}
\usepackage{caption}

% --- Code block styling ---
\definecolor{codebg}{RGB}{245,245,245}
\definecolor{codeframe}{RGB}{220,220,220}
\definecolor{codekw}{RGB}{0,0,192}
\definecolor{codecomment}{RGB}{0,128,0}
\definecolor{codestring}{RGB}{163,21,21}
\definecolor{codenum}{RGB}{128,0,128}
\definecolor{codepunct}{RGB}{90,90,90}
\definecolor{badgegreen}{RGB}{34,139,34}
\definecolor{badgeblue}{RGB}{30,144,255}
\definecolor{badgeblack}{RGB}{50,50,50}

% --- Difficulty badges (rounded corners via tikz, pre-rendered in saveboxes) ---
\usepackage{tikz}
\newsavebox{\badgechubox}
\savebox{\badgechubox}{\tikz[baseline]{\node[fill=badgegreen,text=white,rounded corners=2.5pt,inner xsep=4pt,inner ysep=1.5pt,font=\sffamily\footnotesize\bfseries]{初级};}}
\newsavebox{\badgezhongbox}
\savebox{\badgezhongbox}{\tikz[baseline]{\node[fill=badgeblue,text=white,rounded corners=2.5pt,inner xsep=4pt,inner ysep=1.5pt,font=\sffamily\footnotesize\bfseries]{中级};}}
\newsavebox{\badgegaobox}
\savebox{\badgegaobox}{\tikz[baseline]{\node[fill=badgeblack,text=white,rounded corners=2.5pt,inner xsep=4pt,inner ysep=1.5pt,font=\sffamily\footnotesize\bfseries]{高级};}}
\newcommand{\lvchu}{\usebox{\badgechubox}}
\newcommand{\lvzhong}{\usebox{\badgezhongbox}}
\newcommand{\lvgao}{\usebox{\badgegaobox}}
\pdfstringdefDisableCommands{%
  \def\lvchu{[初级]}%
  \def\lvzhong{[中级]}%
  \def\lvgao{[高级]}%
}

\lstset{
  basicstyle=\ttfamily\scriptsize,
  keywordstyle=\color{codekw}\bfseries,
  commentstyle=\color{codecomment}\itshape,
  stringstyle=\color{codestring},
  backgroundcolor=\color{codebg},
  frame=single,
  rulecolor=\color{codeframe},
  breaklines=true,
  breakatwhitespace=false,
  breakindent=0pt,
  postbreak=\mbox{\textcolor{red}{$\hookrightarrow$}\space},
  numbers=left,
  numberstyle=\tiny\color{gray},
  columns=flexible,
  keepspaces=true,
  showstringspaces=false,
  tabsize=2,
  aboveskip=4pt,
  belowskip=4pt,
  extendedchars=true,
  literate={，}{{，}}1 {；}{{；}}1 {！}{{！}}1 {？}{{？}}1 {“}{{“}}1 {”}{{”}}1 {（}{{（}}1 {）}{{）}}1 {《}{{《}}1 {》}{{》}}1,
}

% --- Extra language definitions (no built-in listings support) ---
\lstdefinelanguage{json}{
  morekeywords={true,false,null},
  sensitive=true,
  morestring=[b]",
  comment=[l]{//},
  morecomment=[s]{/*}{*/},
}
\lstdefinelanguage{yaml}{
  morekeywords={true,false,null,yes,no,on,off},
  sensitive=false,
  morecomment=[l]{\#},
  morestring=[b]",
  morestring=[d]',
}
\lstdefinelanguage{http}{
  morekeywords={GET,POST,PUT,DELETE,PATCH,HEAD,OPTIONS,HTTP,Host,Accept,Authorization,Content-Type,Content-Length,User-Agent},
  sensitive=true,
  morecomment=[l]{\#},
  morestring=[b]",
}
\lstdefinelanguage{TypeScript}{
  morekeywords={abstract,any,as,async,await,break,case,catch,class,const,continue,debugger,declare,default,delete,do,else,enum,export,extends,false,finally,for,from,function,get,if,implements,import,in,instanceof,interface,keyof,let,module,namespace,never,new,null,number,of,package,private,protected,public,readonly,return,set,static,string,super,switch,symbol,this,throw,true,try,type,typeof,undefined,unknown,var,void,while,with,yield,Record,Array,Promise,AsyncIterable,ReadableStream,Date,Error,Object,JSON},
  sensitive=true,
  morecomment=[l]{//},
  morecomment=[s]{/*}{*/},
  morestring=[b]",
  morestring=[b]',
  morestring=[b]`{`},
}

% --- Compact spacing ---
\linespread{0.95}
\setlength{\parskip}{1pt plus 1pt minus 1pt}
\setlength{\parindent}{0pt}

% --- Table styling ---
\renewcommand{\arraystretch}{1.1}

% --- Heading styling (numbered) ---
\titleformat{\section}{\normalsize\bfseries}{\thesection\quad}{0em}{}
\titlespacing{\section}{0pt}{6pt}{2pt}
\titleformat{\subsection}{\small\bfseries}{\thesubsection\quad}{0em}{}
\titlespacing{\subsection}{0pt}{4pt}{1pt}
\titleformat{\subsubsection}{\small\bfseries}{\thesubsubsection\quad}{0em}{}
\titlespacing{\subsubsection}{0pt}{3pt}{1pt}

% --- Page style ---
\pagestyle{fancy}
\fancyhf{}
\fancyhead[L]{\small\emph{\leftmark}}
\fancyhead[R]{\small\thepage}
\renewcommand{\headrulewidth}{0.4pt}
\renewcommand{\sectionmark}[1]{}  % prevent sections from overwriting leftmark

% --- Hyperref setup ---
\hypersetup{
  colorlinks=true,
  linkcolor=blue,
  urlcolor=blue,
  citecolor=blue,
}

% --- Caption format ---
\DeclareCaptionFormat{myformat}{\small\bfseries #1#2#3}
\captionsetup{format=myformat}

\begin{document}

\title{Agent 技术雷达}
\date{}
\maketitle
\markboth{Java 版 Agent 知识}{ Java 版 Agent 知识}

\begin{quote}
语言策略：\texttt{shared}。三语言内容一致，Java 为 canonical。 地图节点：\texttt{08 生态与前沿 / 技术雷达}。本雷达为知识库基线（2026.08），季度更新；采用 Adopt / Trial / Assess / Hold 四级。
\end{quote}



\section*{0. 使用说明}

\begin{itemize}[itemsep=1pt, topsep=2pt]
  \item \textbf{Adopt}：知识库已有实践，默认采用；
  \item \textbf{Trial}：值得试点，但有前提；
  \item \textbf{Assess}：值得观察，暂不投入；
  \item \textbf{Hold}：当前不建议生产采用。
\end{itemize}


\section*{1. 协议与标准}

{\footnotesize
\begin{tabular}{|p{\dimexpr 0.396\columnwidth-2\tabcolsep\relax}|p{\dimexpr 0.125\columnwidth-2\tabcolsep\relax}|p{\dimexpr 0.479\columnwidth-2\tabcolsep\relax}|}
\hline
\textbf{项目} & \textbf{级别} & \textbf{说明} \\
\hline
MCP & Adopt & 工具/资源/提示词接入标准，知识库已纳入工具层 \\
\hline
Agent Skills & Adopt & 经验包与延迟加载模式，适合团队知识沉淀 \\
\hline
A2A & Assess & 多智能体任务交换方向正确，生态仍在收敛 \\
\hline
OpenTelemetry GenAI & Trial & Trace 语义正在稳定，可用于供应商无关采集 \\
\hline
\end{tabular}
}



\section*{2. 框架与库}

{\footnotesize
\begin{tabular}{|p{\dimexpr 0.411\columnwidth-2\tabcolsep\relax}|p{\dimexpr 0.107\columnwidth-2\tabcolsep\relax}|p{\dimexpr 0.482\columnwidth-2\tabcolsep\relax}|}
\hline
\textbf{项目} & \textbf{级别} & \textbf{说明} \\
\hline
LangGraph & Adopt & Python 生产级 Graph + Agent 基线 \\
\hline
LangGraph.js & Adopt & TypeScript 生产级基线 \\
\hline
Spring AI Alibaba Graph & Adopt & Java 生态基线 \\
\hline
OpenAI Agents SDK & Trial & 原型快，注意供应商绑定 \\
\hline
AutoGen & Assess & 多智能体研究价值高，生产运维需自建 \\
\hline
CrewAI & Assess & 角色化流水线友好，复杂控制流要评估 \\
\hline
\end{tabular}
}



\section*{3. 平台与产品}

{\footnotesize
\begin{tabular}{|p{\dimexpr 0.483\columnwidth-2\tabcolsep\relax}|p{\dimexpr 0.100\columnwidth-2\tabcolsep\relax}|p{\dimexpr 0.417\columnwidth-2\tabcolsep\relax}|}
\hline
\textbf{项目} & \textbf{级别} & \textbf{说明} \\
\hline
Claude Code / Codex 类编码 Agent & Adopt & 已有 Harness/Loop 实践素材 \\
\hline
LangSmith / Langfuse 类观测平台 & Trial & 与自建 Trace 二选一，先定义 Span 模型 \\
\hline
通用 Agent 托管平台 & Assess & 关注锁定、数据边界与单位成本 \\
\hline
低代码 Agent 平台 & Hold & 复杂控制与可观测性不足时不要硬上 \\
\hline
\end{tabular}
}



\section*{4. 研究趋势}

{\footnotesize
\begin{tabular}{|p{\dimexpr 0.490\columnwidth-2\tabcolsep\relax}|p{\dimexpr 0.122\columnwidth-2\tabcolsep\relax}|p{\dimexpr 0.388\columnwidth-2\tabcolsep\relax}|}
\hline
\textbf{趋势} & \textbf{级别} & \textbf{对工程的影响} \\
\hline
Agentic Workflows & Adopt & 固定骨架 + 局部自主是当前生产主流 \\
\hline
Context Engineering & Adopt & 决定效果上限，优先于换模型 \\
\hline
Memory 分层 & Trial & 短期上下文 + 长期检索 + 项目记忆 \\
\hline
Multi-Agent & Assess & 只在单 Agent 证明不足后采用 \\
\hline
Agent 评测与安全 & Adopt & 没有回归门禁与护栏不允许上线 \\
\hline
GUI / Computer Use Agent & Assess & 观察沙箱与审计成熟度 \\
\hline
\end{tabular}
}



\section*{5. 更新规则}

\begin{enumerate}[itemsep=1pt, topsep=2pt]
  \item 季度评审一次，更新级别和说明；
  \item 每个项目必须有对应文档或实践记录，不允许只贴链接；
  \item 从 Assess 升到 Trial 需要 PoC 记录；升到 Adopt 需要评测与回滚方案。
\end{enumerate}


\section*{6. 延伸阅读}

\begin{itemize}[itemsep=1pt, topsep=2pt]
  \item \href{../04-工程实践/07-Agent框架选型对比.md}{Agent 框架与平台横评}
  \item \href{02-论文精读清单.md}{论文精读清单}
  \item \href{../../../知识地图总索引.md}{知识地图总索引}
\end{itemize}

\end{document}
